% Canonical mathematical source. Compile outside the repository tree.
\documentclass[11pt]{article}
\usepackage[TU]{fontenc}
\usepackage{lmodern}
\usepackage{amsmath,amssymb,amsthm}
\usepackage[margin=1in]{geometry}
\usepackage{microtype}
\usepackage[colorlinks=true,linkcolor=blue,citecolor=blue,urlcolor=blue]{hyperref}
\hypersetup{pdftitle={Sude's prime-sum squares},pdfsubject={A construction step by step},pdfauthor={Mira}}
\setlength{\emergencystretch}{2em}
\newtheorem{theorem}{Theorem}
\newtheorem{lemma}[theorem]{Lemma}
\newtheorem{proposition}[theorem]{Proposition}
\newtheorem{corollary}[theorem]{Corollary}
\theoremstyle{definition}
\newtheorem{example}[theorem]{Example}
\title{Sude's prime-sum squares\\[0.3em]\large A construction step by step}
\author{Mira}
\date{12 November 2026}
\begin{document}
\maketitle

\begin{abstract}
Sude's square problem asks whether the first \(n^2\) odd primes can be arranged
in an \(n\times n\) square so that every row, every column and both main diagonals
have prime sums, with each prime used exactly once. We follow the large-order
construction through three stages. First, choose prime targets for the row
and column sums with the correct total: nearby prime choices leave a small
error, which a reserved set of choices corrects exactly. Second, place the
prescribed entries, reserving pairs whose swaps transfer amounts between line
sums; an integer-flow argument selects swaps that attain the targets. Third,
use reserved rectangle swaps to make both diagonal sums equal to a common
prime without changing any row or column sum. We prove the combinatorial steps
as they arise, explain the bounds they require, and state the quantitative
prime-supply and placement results used to meet those bounds. Together with
exact certificates for the 1,000 odd orders below 2000, the cited construction
results give the all-odd conclusion under the specified analytic and finite inputs.
\end{abstract}

\section{The square we want to build}\label{sec:problem}

Let \(n\) be a positive integer, the number of rows and columns of the square;
we call \(n\) its \emph{order}. Let \(p_j\) be the \(j\)th prime, with
\(p_1=2\), and set
\[
 P_n=\{p_2,p_3,\ldots,p_{n^2+1}\}.
\]
An arrangement is a bijection
\(A:\{0,\ldots,n-1\}^2\longrightarrow P_n\). Write
\[
 R_i(A)=\sum_{j=0}^{n-1}A(i,j),\qquad
 C_j(A)=\sum_{i=0}^{n-1}A(i,j)
\]
for its row and column sums, and
\[
 D_+(A)=\sum_{i=0}^{n-1}A(i,i),\qquad
 D_-(A)=\sum_{i=0}^{n-1}A(i,n-1-i)
\]
for its main-diagonal and antidiagonal sums. We require all \(2n+2\) of these
labelled sums to be prime. They may repeat. At \(n=1\), all four labels refer
to the entry 3.

Odd order is necessary: when \(n\ge2\) is even, every row sum is even and
at least \(3n>2\). At order 3 there is, for example, the square
\[
 \begin{pmatrix}7&5&19\\13&11&29\\17&3&23\end{pmatrix}.
\]
Its row sums are \(31,53,43\), its column sums are \(37,19,71\), and its diagonal
sums are \(41,47\). These use exactly the first nine odd primes.

The construction below concerns a fixed odd order \(n\ge2001\). We first choose
the prime sums we want. We then prepare one arrangement with enough reserved
swaps to attain those sums, and finally use a different family of swaps for
the diagonals. All estimates refer to this same arrangement. The smaller odd
orders are handled by the finite certificates in Section~\ref{sec:allorders}.

We prove the elementary selection, switching and routing arguments in this
paper. The large prime-supply estimates and the joint placement and completion
bounds are used as quantitative results of the supporting construction
\cite[Sections 3--4]{Companion}, with the published analytic inputs collected
in Appendix~\ref{app:inputs}. Their role is stated where each is needed.
Thus the walkthrough explains how the supplied objects produce a square,
without treating the long quantitative estimates as consequences of the
short combinatorial lemmas alone.

\section{First choose the prime sums}\label{sec:targets}

Let
\[
 T=\sum_{p\in P_n}p,\qquad X=\max P_n,\qquad y=T/n.
\]
Every arrangement has total row sum \(T\) and total column sum \(T\). Thus \(y\)
is the common average of the row sums and of the column sums; the average
\emph{entry} is \(T/n^2\). We aim to choose odd primes \(t_0,\ldots,t_{n-1}\)
with
\begin{equation}\label{eq:targets}
 \sum_i t_i=T,\qquad |t_i-y|\le D,
\end{equation}
where \(D\ge0\) is a target radius. We will use the same vector for the row and
column sums. Target values may repeat and need not belong to \(P_n\): they
are sums of entries, not entries to be placed.

The exact total is indispensable. For the nine-entry pool above, \(T=127\),
so choosing three target sums of 43 cannot work, although each target is prime.
The triple \((41,43,43)\) has the right total. That observation supplies targets,
not yet a placement attaining them.

\subsection{Choose near the mean, then correct the total}

Choosing primes near \(y\) independently need not give total \(T\). Instead,
reserve some target positions for a final exact correction. The other positions
can be filled greedily without allowing their total error to accumulate.

\begin{lemma}[Greedy prime brackets]\label{lem:brackets}
Suppose each of \(m\) positions has two permitted values \(\ell_i,u_i\) with
\(y-D\le\ell_i\le y\le u_i\le y+D\). One can choose one value at each position
so that their sum \(V\) satisfies \(|V-my|\le D\).
If \(\ell_i\equiv u_i\pmod h\) at every position, the chosen sum also satisfies
\(V\equiv\sum_i\ell_i\pmod h\).
\end{lemma}
\begin{proof}
Start with discrepancy zero. When the discrepancy of the choices already made
from their prescribed mean is nonnegative, choose the next lower value;
otherwise choose the upper value. If the old discrepancy lies in \([0,D]\),
adding \(\ell_i-y\in[-D,0]\) keeps it in \([-D,D]\). If it lies in \([-D,0)\),
adding \(u_i-y\in[0,D]\) does the same. Induction proves the error bound. The
congruence is unchanged by either choice at each position.
\end{proof}

At each of the remaining \(s=n-m\) positions, suppose we have two odd primes
\(a_j,a_j+g_j\) in \([y-D,y+D]\), where \(g_j\ge0\). Initially choose every
\(a_j\), and put \(L=\sum_j a_j\). Replacing selected lower choices by their
upper choices adds a subset sum of the gaps \(g_j\).

Both groups of positions offer a lower and an upper prime. Their roles differ:
the first \(m\) choices are made greedily to keep their sum \(V\) close to
\(my\), while the other \(s\) choices are reserved until \(V\) is known and
then used to correct the total exactly. The reserved pairs need not individually
straddle \(y\); it is their combined gap coverage that matters.

\begin{lemma}[Exact total from a correction interval]\label{lem:exact-targets}
Let \(T=ny\) be an integer and \(m+s=n\). Suppose the \(m\) positions of
Lemma~\ref{lem:brackets} offer odd primes in \([y-D,y+D]\), and let \(V\) be
the sum selected by its greedy procedure. Suppose the \(s\) reserved pairs
\(a_j,a_j+g_j\) also consist of odd primes in that interval. Let \(h\) be a
positive integer. Assume the reserved gaps \(g_j\) have subset sums containing every
multiple of \(h\) in an integer interval \([a,b]\), where \(0\le a\le b\), and
\begin{equation}\label{eq:target-margin}
 a+D\le sy-L,\qquad sy-L+D\le b.
\end{equation}
If \(V\equiv T-L\pmod h\), one can select a prime from each reserved pair so
that, together with the \(m\) greedily selected primes, they give \(n\) targets
satisfying \eqref{eq:targets}.
\end{lemma}
\begin{proof}
The amount still needed is \(r=T-V-L\). Since
\[
 r=sy-L-(V-my),
\]
Lemma~\ref{lem:brackets} and \eqref{eq:target-margin} give \(a\le r\le b\).
The congruence gives \(h\mid r\). Choose a subset \(J\) with
\(\sum_{j\in J}g_j=r\), and use the upper prime precisely at those positions.
The total is now \(V+L+r=T\); every chosen value is still an odd prime within
\(D\) of \(y\).
\end{proof}

In the simplest case \(h=2\), the required congruence is automatic. The total
\(T\) is a sum of \(n^2\) odd numbers, so \(T\equiv n\pmod2\), whereas
\(V+L\equiv m+s=n\pmod2\). The residual is therefore even. At \(n=3\),
\(y=127/3\), take \(D=4/3\), one bracket \((41,43)\) for the greedy step,
and two copies of \((41,43)\) reserved for correction. The greedy step chooses
41. The residual above the three lower choices is \(127-123=4\); choosing 43
from both reserved pairs gives \((41,43,43)\). Repeating a prime here is allowed
because these are target sums.

For larger even \(h\), prime brackets in selected reduced residue classes
can enforce the congruence. The following elementary fact explains the choice
of classes; the existence of primes in those classes is a separate analytic input.

\begin{lemma}[Choosing residue classes for the greedy step]\label{lem:units}
If \(h\ge2\) is even, \(m\ge2\), and \(c\equiv m\pmod2\), then \(c\) is a sum of
\(m\) residues coprime to \(h\), modulo \(h\).
\end{lemma}
\begin{proof}
Use \(m-2\) copies of 1. The residual \(d=c-(m-2)\) is even. For the power of 2
dividing \(h\), choose an odd residue \(a\), so \(d-a\) is odd as well. For
each odd prime divisor \(p\) of \(h\), choose \(a\) modulo \(p\) different from
both 0 and \(d\); at most two residues are excluded and \(p\ge3\). Lift these
choices to the corresponding prime powers and combine them by the Chinese
remainder theorem. Both \(a\) and \(d-a\) are coprime to \(h\). They provide
the last two residues.
\end{proof}

Apply this with \(c=T-L\), whose parity is \(n-s=m\). For each chosen
unit class, supply a lower and upper prime bracket in that same class.
Lemma~\ref{lem:brackets} then preserves the required total residue while
keeping the greedily chosen sum within \(D\) of \(my\).

\subsection{Build a correction interval rather than assume one}

We next explain the two requirements in Lemma~\ref{lem:exact-targets}: gap
coverage and enough room around the desired residual.

\begin{lemma}[Extending subset-sum coverage]\label{lem:grid}
Suppose subset sums of some nonnegative gaps include
\(0,h,2h,\ldots,R\), where \(h>0\) and \(h\mid R\). Adjoining a gap \(g\)
with \(h\mid g\) and \(0\le g\le R+h\) extends this coverage through \(R+g\).
\end{lemma}
\begin{proof}
The old choices give multiples of \(h\) in \([0,R]\), and choices using the
new gap give them in \([g,R+g]\). The two lists overlap or are adjacent in
steps of \(h\), so leave no missing multiple.
\end{proof}

Copies of a seed gap \(h\) start such a list; suitable further gaps extend it.
For example, gaps 2, 4 and 6 realize every even integer from 0 through 12.
Gaps 4 and 8 have the same total but miss 2, 6 and 10. A total amount of
available change is therefore not a substitute for complete coverage.
When \(a\) and \(b\) are multiples of \(h\), adjoining a nonnegative gap
\(g\) divisible by \(h\) similarly extends coverage from \([a,b]\) to
\([a,b+g]\), provided \(g\le b-a+h\).

Here is a quantitative way to obtain the margins when the correction interval
is \([0,G]\), with \(G=\sum_jg_j\). Put \(M=L+G/2\), the sum of the pair
midpoints. The two inequalities \eqref{eq:target-margin} are equivalent to
\begin{equation}\label{eq:midpoint-margin}
 |M-sy|+D\le G/2.
\end{equation}
Pairs with large gaps can supply this room while keeping their midpoints near
\(y\).

\begin{lemma}[A midpoint margin]\label{lem:margin}
Suppose the positions reserved for correction consist of \(P\) pairs with midpoints within
\(D/2\) of \(y\), and \(W\) further pairs with midpoints within \(D/4\) of
\(y\) and gaps at least \(D\). If \(W\ge2P+5\), then
\[
 |M-sy|+D+D/4\le G/2.
\]
In particular, \eqref{eq:midpoint-margin} holds.
\end{lemma}
\begin{proof}
The triangle inequality gives \(|M-sy|\le PD/2+WD/4\), and the wide pairs give
\(G/2\ge WD/2\). The difference between these two bounds is
\((W-2P)D/4\ge5D/4\).
\end{proof}

The prime-target supply results of \cite[Section 3.2 and the range constructions
in Section 4]{Companion} implement these steps: obtain prime brackets, construct
gap coverage, verify a midpoint margin, and check that the greedy and reserved
positions fit among the \(n\) targets. In the bounded lower range they use
short seed gaps and prime brackets in a finite set of residue classes. In the
unbounded range additional gap choices cover the even residue classes of a
larger seed step, producing the correction interval required by
Lemma~\ref{lem:exact-targets}. The elementary lemmas above explain why these
supplies give an \emph{exact} total, rather than merely a small error.

For later use, the lower-range target radii are
\begin{equation}\label{eq:radii}
 D=\begin{cases}
 (23/50)X,&2001\le n\le9299,\\
 y/400,&9301\le n\le19483,
 \end{cases}
\end{equation}
and the unbounded construction uses \(D=y/1400\) for \(n\ge35735\).
The middle range has its own target and capacity estimates in the same cited
construction. We will now see exactly where a bound on \(D\) is needed.

\section{Reserve the swaps before filling the square}\label{sec:reserve}

The targets have been chosen. We must now arrange the \emph{distinct entries}
of \(P_n\) so that enough adjustments remain available. Unlike target choices,
physical pairs cannot reuse an entry or a cell.

\subsection{A row transfer we can perform}

If entries \(a,b\) occupy two cells in the same column, swapping them changes
the two row sums by \(b-a\) and \(a-b\). All column sums stay fixed. Since
both entries are odd, the changes are even. If both cells avoid the two
diagonals, their sums stay fixed too. A corresponding swap within a row
transfers an even amount between two column sums.

Call a collection of reserved, disjoint pairs of this kind a \emph{reservoir}.
Its possible adjustments are the sums of the changes supplied by the pairs
we choose to swap. For example, consider these six entries in two rows:
\[
 \begin{pmatrix}3&7&19\\5&11&13\end{pmatrix}.
\]
Swapping within its three columns changes the first row sum by a subset sum
of \(2,4,-6\), namely any of
\[
 -6,-4,-2,0,2,4,6.
\]
The second row changes by the negative of that amount and every column sum
is unchanged. In units of 2, this reservoir has signed capacity 3. This is a
local example: its cells must still be placed away from the diagonals and all
other reserved objects in the full square.

More generally, suppose a reservoir's pair gaps are \(2w_1,\ldots,2w_k\), and
subset sums of the \(w_j\) contain every integer from 0 through \(W\). Begin
with the smaller entries in one row, and swap a subset of total weight
\(K=\lfloor W/2\rfloor\). From this \emph{centered orientation}, choosing any
other subset allows changes \(2(z-K)\), \(0\le z\le W\). In particular it
allows every even change between \(-2K\) and \(2K\). This explains both the
centering step and the division by 2 in the eventual flow calculation.

\subsection{Reserve a different kind of swap for the diagonals}\label{subsec:rectangles}

A pair swap generally changes row or column sums. To adjust a diagonal after
those sums are finished, reserve four entries at the corners of a rectangle:
\begin{equation}\label{eq:trade}
 \begin{pmatrix}a&a+d\\b+d&b\end{pmatrix}
 \quad\longleftrightarrow\quad
 \begin{pmatrix}a+d&a\\b&b+d\end{pmatrix}.
\end{equation}
The four values are distinct odd primes from \(P_n\), and \(d>0\) is even.
The rows or columns of the rectangle need not be adjacent.

\begin{lemma}[Rectangle switch]\label{lem:rectangle}
The switch \eqref{eq:trade} preserves the entries and every row and column sum.
For any set of cells \(L\), its sum changes by
\[
 d\bigl(\mathbf1_L(r_1,c_1)-\mathbf1_L(r_1,c_2)
       -\mathbf1_L(r_2,c_1)+\mathbf1_L(r_2,c_2)\bigr),
\]
where \(\mathbf1_L(r,c)\) is 1 if \((r,c)\in L\), and 0 otherwise.
\end{lemma}
\begin{proof}
The changes at the four corners, in row-major order, are \(d,-d,-d,d\).
They cancel within each affected row and column. The displayed formula sums
exactly the changes belonging to \(L\).
\end{proof}

For example,
\[
 \begin{pmatrix}3&5\\13&11\end{pmatrix}
 \longleftrightarrow
 \begin{pmatrix}5&3\\11&13\end{pmatrix}
\]
keeps row sums 8 and 24 and column sums 16 and 16. At order 5, place its
corners in rows 0,1 and columns 0,2. Only its upper-left corner lies on the
main diagonal, and none lies on the antidiagonal. The switch therefore
changes the diagonal pair by \((2,0)\). Reflecting this placement in the
vertical axis gives the analogous effect on the antidiagonal. These illustrate
separate placements; a family of such objects must be packed without common cells.

For the full construction, reserve one family whose forward switches change
only \(D_+\), and a second family whose forward switches change only \(D_-\).
All rectangle supports are disjoint. Their even gaps give complete subset-sum
intervals. Relative to the chosen orientations, let the main-diagonal family
allow every even change in an integer interval \([\ell_+,u_+]\), and let the
antidiagonal family allow every even change in \([\ell_-,u_-]\). The intervals
may consist of positive changes or may extend on both sides of zero. We retain
the actual endpoints: replacing an interval by a smaller symmetric one could
lose a needed target. The rectangles also avoid the central cell, so the common
cell of the diagonals creates no unintended interaction between their switches.

\subsection{Choose entries without using any twice}

The actual supply question is now visible. We need prime pairs with useful
gaps, four distinct entries with the equal-gap pattern \eqref{eq:trade}, and
unused entries for all remaining cells. The selected pair and rectangle
families must be disjoint from one another.

Here is the elementary counting principle used when excluding earlier choices.
If \(\mathcal F\) is a finite family of candidate objects and every already
used entry lies in at most \(d\) members of \(\mathcal F\), then a used set
\(U\) excludes at most \(d|U|\) candidates. Consequently,
\(|\mathcal F|>d|U|\) guarantees an available candidate. This follows by counting
all incidences between used entries and candidates; counting an excluded
candidate more than once only enlarges the upper bound. More complicated
selection steps use a sum of such exclusion bounds.

To see how analytic estimates provide candidates, define
\[
 \theta(x;q,a)=\sum_{\substack{p\le x\\p\equiv a\pmod q}}\log p,
\]
where the sum is over primes. Let \(\varphi(q)\) denote Euler's totient function.
Suppose \(q\ge1\), \(2\le u<v\), \(\gcd(a,q)=1\), and available estimates give
\[
 |\theta(u;q,a)-u/\varphi(q)|\le E_u,\qquad
 |\theta(v;q,a)-v/\varphi(q)|\le E_v.
\]
Subtracting shows that the number of primes in \((u,v]\) in this class is at least
\begin{equation}\label{eq:stock}
 \frac{(v-u)/\varphi(q)-E_u-E_v}{\log v},
\end{equation}
because each prime counted contributes at most \(\log v\). In particular,
a positive numerator supplies an odd prime. If the numerator exceeds
\(F\log v\), there are more than \(F\) such primes, so deleting at most
\(F\) of them still leaves one. For entries of the square we also require
\(v\le X\); for prime \emph{target sums} that restriction is unnecessary.

Prime density also starts gap selection: \(k\ge2\) primes in an interval of
length \(H\) have two consecutive members at distance at most \(H/(k-1)\),
since their \(k-1\) gaps sum to at most \(H\). Further selection and congruence
arguments turn such gaps into the complete adjustment sets above. The
published estimates used for these counts and for bounding excluded
progressions are listed in Appendix~\ref{app:inputs}. The quantitative
construction checks their domains and the remaining counts after each selection.

\section{Complete one reference arrangement}\label{sec:reference}

\subsection{Place the pairs without collisions}

There is a useful elementary answer to the placement question for one layer
of row pairs. Let \(\sigma\) be a permutation of the rows with no fixed points.
For each row label \(i\), one pair is to connect rows \(i\) and \(\sigma(i)\).
We must choose a column for that pair, avoiding reserved cells and both diagonals.
Using a different column for each pair also prevents collisions within the layer.

\begin{lemma}[Placing a layer of pairs]\label{lem:layer}
Suppose the already reserved cells occupy at most \(L\) cells in each row and
each column. If \(4L+8<n\), there is a bijection \(f\) from the pair labels
to the columns such that the cells \((i,f(i))\) and \((\sigma(i),f(i))\) avoid
all reservations and both diagonals. The resulting pairs have disjoint cells,
and add exactly two occupied cells to each row and each column.
\end{lemma}
\begin{proof}
For pair \(i\), at most \(2L\) columns are excluded by reservations in its two
rows. At most four more are excluded by the two diagonals. Conversely, a
fixed column excludes at most \(2L\) pair labels because of reservations:
at most \(L\) whose first row is occupied, and at most \(L\) whose second row
is occupied, since \(\sigma\) is a permutation. Diagonal avoidance excludes
at most four additional labels.

Form a bipartite graph of pairs and columns, joining precisely the allowed
choices. Each vertex has at most \(b=2L+4\) forbidden partners, with \(2b<n\).
For any nonempty set \(S\) of pair labels with \(|S|\le n/2\), even one of its
members has at least \(n-b>n/2\) allowed columns, so the union of their
neighbors has size at least \(|S|\). If \(|S|>n/2\), every column has a neighbor
in \(S\), since otherwise it would forbid more than \(b\) labels. Thus the
neighbor condition holds for every \(S\), including the empty set.

This condition gives a matching of all pair labels: if a maximum matching
left a label unmatched, follow alternating unmatched and matched edges from
it. Unless an unmatched column is reached, giving a larger matching, the
reachable labels have fewer neighboring columns than labels. That contradicts
the condition just proved. The matching is the required bijection \(f\).
Distinct columns separate different pairs. Each row is the first row of one
pair and the second row of one other pair, and these use different columns
because \(\sigma\) has no fixed point.
\end{proof}

This is the placement bound used for permutation layers in the supporting
construction \cite[Section 3.1]{Companion}. Repeating it installs further layers
as long as the updated reservation count satisfies the inequality. Transposing
the square gives the corresponding statement for column pairs, with all
previous row pairs now counted among the reservations. The quantitative
construction chooses the number of layers, their row permutations and their
prime gaps to meet both this placement bound and the capacity bounds below.

\subsection{Fill the remaining cells and measure the same square}

We now install the selected primes in their designated cells. Here
\emph{install} means place the pairs and rectangles with exactly the switching
effects just proved. Row reservoirs, column reservoirs and diagonal rectangles
have pairwise disjoint supports. Ordinary row and column reservoirs avoid
both diagonals. The ordinary row and column reservoirs are centered before
we measure discrepancies. The diagonal families retain the specified
orientations and adjustment intervals of Section~\ref{subsec:rectangles}.

If \(k\) distinct pool entries occupy \(k\) distinct cells, the remaining
\(n^2-k\) entries can always be bijected with the remaining cells. This proves
that the partial placement can be completed without changing the pool. But
an arbitrary completion could leave row sums too far from their targets. The
quantitative placement and completion results choose a completion \(A_0\)
whose row and column sums obey
\begin{equation}\label{eq:baseline}
 |R_i(A_0)-y|\le E,\qquad |C_j(A_0)-y|\le E
\end{equation}
for a controlled error \(E\), while retaining all reserved objects. They also
bound its two diagonal sums and the adjustment intervals available there.
This is one simultaneous placement statement, not separate choices of a good
square for each family of lines.

One elementary device converts a bound on the spread of line sums into
\eqref{eq:baseline}.
\begin{lemma}[From width to distance from the mean]\label{lem:width}
If \(n\ge1\) and \(z_1,\ldots,z_n\) have average \(y\) and \(|z_i-z_j|\le E\) for every
\(i,j\), then \(|z_i-y|\le E\) for every \(i\).
\end{lemma}
\begin{proof}
Write \(z_i-y=n^{-1}\sum_j(z_i-z_j)\) and use the triangle inequality.
\end{proof}

This applies separately to row sums and column sums because both totals are
\(T\). The substantial completion estimates control their widths despite the
already reserved cells. The relevant joint installation and completion
arguments are part of the lower, middle and unbounded constructions cited
in Section~\ref{sec:ranges}; their conclusion supplies the single \(A_0\)
used from this point onward.

\section{Route the row and column discrepancies}\label{sec:routing}

For each row, define the amount it must \emph{give away}, in units of 2, by
\[
 b_i^R=\frac{R_i(A_0)-t_i}{2}.
\]
A negative value means the row must receive that amount instead. These demands
are integers: an odd number of odd entries has odd sum, as does the target.
They sum to zero because \(\sum_iR_i(A_0)=\sum_it_i=T\). Define \(b_j^C\)
for the columns in the same way. From \eqref{eq:targets} and
\eqref{eq:baseline},
\begin{equation}\label{eq:demand-bound}
 |b_i^R|,|b_j^C|\le (E+D)/2.
\end{equation}
This is the promised use of the target radius: the closer the targets and
baseline sums are to the same mean, the less switching capacity we need.

Draw a graph whose vertices are rows. An edge records a reserved reservoir
between its two rows. Fix an orientation for each edge. A signed integer
\(f_e\) means transferring \(2f_e\) from its initial row to its terminal row;
a negative value reverses the transfer. If the reservoir can realize every
such transfer for \(|f_e|\le c_e\), call \(c_e\) its capacity. Disjoint supports
allow edge choices to be performed together. Thus we seek a flow satisfying
\(\operatorname{div}f(i)=b_i^R\), where divergence is outgoing minus incoming
flow. The net row change is then \(-2b_i^R=t_i-R_i(A_0)\).

\begin{lemma}[Integral routing criterion]\label{lem:flow}
Let a finite graph have integer capacities \(c_e\ge0\) and integer demands
\(b_v\) with total zero. A signed integer flow with
\(\operatorname{div}f(v)=b_v\) and \(|f_e|\le c_e\) exists if and only if,
for every vertex set \(S\),
\begin{equation}\label{eq:cut}
 \left|\sum_{v\in S}b_v\right|\le\sum_{e\in\delta(S)}c_e.
\end{equation}
Here \(\delta(S)\) is the set of edges crossing from \(S\) to its complement.
\end{lemma}
\begin{proof}
Summing divergences over \(S\) cancels every internal edge, proving necessity.
For sufficiency, replace each edge by two opposite arcs, each of capacity
\(c_e\). Add a source with capacity \(b_v\) into each vertex of positive demand,
and an arc of capacity \(-b_v\) from each vertex of negative demand to a sink.
Put \(Q=\sum_{b_v>0}b_v\). A source--sink cut whose original vertices on the
source side form \(S\) has capacity
\[
 Q-\sum_{v\in S}b_v+\sum_{e\in\delta(S)}c_e\ge Q.
\]
Starting with zero flow, augment along a residual source--sink path whenever
one exists. Integer capacities permit integer augmentations, each increasing
the value by at least 1. The value is bounded by \(Q\), so the process stops.
If it stopped below \(Q\), the vertices reachable from the source in the
residual graph would give a cut of smaller capacity, a contradiction. All
source and sink arcs are therefore saturated. Subtract the opposite arc flows
on each original edge to obtain the signed flow required.
\end{proof}

\subsection{A readable capacity test}

The cut condition need not be checked separately for every possible demand
vector. Suppose every edge has capacity at least \(C\), every demand has
absolute value at most \(B\), and the graph satisfies
\[
 |\delta(S)|\ge h\min(|S|,n-|S|)
\]
for a constant \(h>0\). Since the total demand is zero, its sum over \(S\)
is bounded using either \(S\) or its complement:
\[
 \left|\sum_{v\in S}b_v\right|
 \le B\min(|S|,n-|S|).
\]
The cut capacity is at least \(Ch\min(|S|,n-|S|)\). Thus
\begin{equation}\label{eq:capacity}
 B\le Ch
\end{equation}
suffices. Combining this with \eqref{eq:demand-bound}, one may choose an
integer \(B\) with \(E+D\le2B\) and then prove \eqref{eq:capacity}.

For example, a specialization used in the lower-range construction has
\(h=199/200\) and \(C=\lfloor N/2\rfloor\), where \(N\ge2000\) is a
lower bound arising from the installed reservoir ranges. Since
\[
 C\ge (N-1)/2\ge 999N/2000,
\]
the numerical condition \(400000B<199\cdot999N\) implies \eqref{eq:capacity}.
The constants express the same simple requirement: the available transfer
across every cut must exceed the amount the rows on that side need to exchange.

\subsection{Apply the two flows to the same square}

First perform the row flow. Each swap preserves all column sums, both
diagonals and the cells reserved for column and diagonal adjustments. We
obtain \(A_1\) with \(R_i(A_1)=t_i\). Its column sums are still those of
\(A_0\), so the originally calculated column demands remain correct.

Now use the column reservoirs to perform the column flow. These swaps
preserve each already finished row sum. The result \(A_2\) therefore satisfies
\[
 R_i(A_2)=t_i,\qquad C_i(A_2)=t_i.
\]
It still agrees with \(A_0\) on the diagonal rectangles, and
\(D_\pm(A_2)=D_\pm(A_0)\). This is why the three sets of reserved objects
were installed together before any routing took place.

\section{Finish both diagonals together}\label{sec:diagonals}

Prime rows and columns do not force prime diagonals. For example, swapping
the first two rows of the small square in Section~\ref{sec:problem} leaves all
row and column sums prime but gives diagonal sums 41 and \(51=3\cdot17\).
We therefore use the rectangles deliberately left untouched in \(A_2\).

Put \(d_+=D_+(A_0)\) and \(d_-=D_-(A_0)\). The two reserved families allow
every even change in \([\ell_+,u_+]\) and \([\ell_-,u_-]\), respectively.
Because their supports are disjoint and each family affects only its own
diagonal, the \emph{joint} set of attainable sums contains
\begin{equation}\label{eq:joint}
 \begin{split}
 &\bigl((d_++2\mathbb Z)\cap[d_++\ell_+,d_++u_+]\bigr)\\
 &\hspace{3em}\times
 \bigl((d_-+2\mathbb Z)\cap[d_-+\ell_-,d_-+u_-]\bigr).
 \end{split}
\end{equation}
Here the product means independent choices of the two coordinates. This
conclusion uses disjoint physical switches, not merely two separate statements
that each diagonal could be adjusted in some arrangement.

\begin{proposition}[A common prime finishes the square]\label{prop:diagonal}
Suppose an odd prime \(q\) belongs to
\begin{equation}\label{eq:overlap}
 [d_++\ell_+,d_++u_+]\cap[d_-+\ell_-,d_-+u_-].
\end{equation}
Then the reserved rectangle switches turn \(A_2\) into a valid prime-sum square.
\end{proposition}
\begin{proof}
Both \(d_+,d_-\) are odd, so \(q-d_+\) and \(q-d_-\) are even. The interval
condition puts them inside the two allowed ranges. Choose the corresponding
switches in each disjoint family. By \eqref{eq:joint}, the choices can be
performed together and make both diagonal sums equal to \(q\). Every
rectangle switch preserves each row and column sum and permutes existing
entries. Hence all required sums are prime and the entries are still exactly
\(P_n\), each once.
\end{proof}

One convenient sufficient condition applies when the actual adjustment
intervals contain symmetric ranges \([-2H_+,2H_+]\) and \([-2H_-,2H_-]\).
If also
\[
 |d_\pm-y|\le U,\qquad U+D\le2H_+,\quad U+D\le2H_-,
\]
then any already chosen target, for example \(q=t_0\), works: the triangle
inequality gives \(|t_0-d_\pm|\le D+U\). No new prime choice is needed under
these additional bounds. They are a useful special case, not a requirement
imposed on every range construction.

In general, retain the actual overlap endpoints
\[
 L=\max(d_++\ell_+,d_-+\ell_-),\qquad
 U'=\min(d_++u_+,d_-+u_-).
\]
The placement bounds must leave an interval \((u,v]\subseteq[L,U']\), with
\(u\ge2\), to which a prime estimate applies. A positive lower bound in
\eqref{eq:stock}, or the corresponding all-scale interval estimate of the
supporting construction, then supplies an odd prime in the overlap. The
lower-range and unbounded constructions use overlap arguments with their
stated interval endpoints. In both versions,
the target is one common prime and the proof of joint reachability is the same.

\section{Where the quantitative bounds enter}\label{sec:ranges}

For our fixed odd \(n\), the walkthrough has reduced the construction to
concrete requirements. We need prime target choices meeting the correction
interval and position budget; distinct pool entries in a joint installation
of row reservoirs, column reservoirs and diagonal rectangles; a completion
with the reference bounds \eqref{eq:baseline}; enough cut capacity for
\eqref{eq:capacity}; and a prime in the diagonal overlap \eqref{eq:overlap}.
The proved steps then act on the single reference arrangement in sequence:
\[
 A_0\ \longrightarrow\ A_1\ \longrightarrow\ A_2\ \longrightarrow\ A_3,
\]
where the arrows finish the rows, then the columns, then both diagonals.

The supporting construction \cite[Sections 3--4]{Companion} establishes these
requirements in three ranges, using the analytic inputs in
Appendix~\ref{app:inputs}. The following are cited quantitative construction
results, not new estimates derived here from prime density alone.

\paragraph{Lower range: odd \(2001\le n\le19483\).}
The proof uses nine certified descriptions of finite prime prefixes and 21
parameter subranges. These give usable prime pairs and explicit bounds on
their gaps. The prime-target selection has radii \eqref{eq:radii}. The joint
placement reserves the pairs and rectangles together, completes the remaining
cells with line-width and diagonal bounds, and verifies the capacity and
overlap inequalities. Its routing step has the numerical specialization
explained in Section~\ref{sec:routing}. Thus all three arrows above apply.
The cited lower-range construction includes these simultaneous placement
bounds, not only a list of targets.

\paragraph{Middle range: odd \(19485\le n\le35733\).}
A verified finite collection of 3,650,551 twin-prime pairs supplies pairs with
the smallest possible nonzero even gap, 2. Further selection stages enlarge
the possible changes while avoiding previously used entries. Prime-supply
lower bounds and short-progression upper bounds control this selection.
The completed square, centered reservoirs, line-error estimates and diagonal
ranges refer to the same placement. The resulting demand and diagonal bounds
permit the same row, column and rectangle operations as above.

\paragraph{Unbounded range: odd \(n\ge35735\).}
The prime-target radius is \(D=y/1400\). As the pool grows, the proof enlarges
the reserved swapping families and controls row and column contributions
simultaneously before completing the unused cells. Its all-scale prime
estimates give fresh pairs and the necessary adjustment ranges, and its
completion estimates retain the diagonal rectangles. The capacity and joint
diagonal arguments then apply at every larger order, not just at a finite
list of tested values.

The exact estimates and range-specific installations reside in the cited
construction source. The proofs in Sections~\ref{sec:targets}--\ref{sec:diagonals}
explain how their bounds are consumed: target radius and reference error
control the demands, reservoir capacity solves them, and diagonal reachability
completes the square. Neither distinctness of the entries nor preservation of
already completed lines is left implicit in this assembly.

\section{Add the finite orders}\label{sec:allorders}

For each odd \(n<2000\), an exact certificate specifies a permutation of the
first \(n^2\) odd primes. A record can be viewed as an initial arrangement
followed by a list of transpositions. Replaying it and summing the resulting
rows, columns and diagonals checks the required square directly.

\paragraph{Finite certificate assertion.}
For every \(n\in\{1,3,\ldots,1999\}\), the stored certificate gives exactly
the prescribed pool and all \(2n+2\) labelled line sums are prime.

The replay checks all 1,000 orders, not a sample. Primality is tested by exact
sieve membership or complete trial division. This is a finite decision
procedure because a composite integer has a prime divisor at most its square
root. An independent replay of the complete family succeeded; the records,
verifier and detailed verification evidence are supplied in
\cite[Sections 6--7]{Companion}. We use that exact computational assertion,
separately from the analytic construction for larger orders.

\begin{theorem}[All positive odd orders]\label{thm:allodd}
Using the quantitative construction results of Section~\ref{sec:ranges} under
the analytic inputs in Appendix~\ref{app:inputs}, and assuming the finite
certificate assertion above, every positive odd integer \(n\) admits a
bijection onto \(P_n\) with prime sums on every row, every column and both
main diagonals.
\end{theorem}
\begin{proof}
The finite assertion covers \(n<2000\). For the other odd orders, use the
lower-range result through 19483, the middle-range result from 19485 through
35733, and the unbounded result from 35735 onward. The integers between
successive endpoints are 2000, 19484 and 35734, all even. Each construction
uses the same pool and the same required lines, so no odd order is omitted.
\end{proof}

The result is an existence construction. It makes no claim of a small running
time at arbitrary order or of optimal thresholds for its separate methods.
Requiring every wrapped diagonal or higher-dimensional line would introduce
additional constraints; those do not follow just by repeating the final
rectangle step. For the present square, the construction succeeds because
it separates the tasks while preserving precisely what the next task needs:
correct total, available disjoint switches, and the completed line sums.

\appendix
\section{Analytic inputs used in the quantitative construction}\label{app:inputs}

For completeness we identify the analytic results used by the cited range
constructions. Write \(\pi(x)\) for the number of primes at most \(x\),
\(\pi(x;q,a)\) for the number congruent to \(a\pmod q\), and
\(\theta(x)=\sum_{p\le x}\log p\). All logarithms are natural. The
progression function \(\theta(x;q,a)\) and Euler's totient \(\varphi(q)\)
were introduced when deriving \eqref{eq:stock}.

\begin{enumerate}
\item \textbf{Pool scale and total.} Rosser and Schoenfeld's prime-count bounds,
Theorem 2 and Corollary 1, equations (3.3)--(3.6); prime-index bounds,
Theorem 3 and its Corollary, equations (3.10)--(3.13); and theta upper bound,
Theorem 4, equation (3.15), supply nine inequalities \cite{RS62}.
They bound the largest available entry \(X\) and the exact-pool mean, so that
intervals chosen for entries remain inside \(P_n\).
\item \textbf{Primes in progressions and intervals.} The theta part of
Theorem 1.9 and equation (1.18), Proposition 2.1 with equations (2.3)--(2.7),
Propositions 2.5 and 2.7 with Definition 2.6, and Proposition 4.34 with
Definition 3.1 of Bennett, Martin, O'Bryant and Rechnitzer supply five inputs
\cite{BMOR18}. The first gives
\[
 |\theta(z;q,a)-z/\varphi(q)|\le2.072\sqrt{x}
 \quad(1\le z\le x\le x_2(q),\ 1\le q\le10^5,\ \gcd(q,a)=1),
\]
with \(x_2(q)\) as defined in that theorem. The others provide the smoothed
explicit formula, zero counts, finite-height zero information and zero-free
region used for the larger intervals. Their stated domains and conductor
conditions are retained. These bounds supply prime brackets, pair candidates
and diagonal targets through interval differences such as \eqref{eq:stock}.
\item \textbf{Upper bounds for excluded progressions.} Montgomery and Vaughan's
Brun--Titchmarsh inequality, Theorem 2, equation (1.10), states for positive
\(x,y\), positive coprime integers \(k,a\), and \(y>k\), that
\[
 \pi(x+y;k,a)-\pi(x;k,a)
 <\frac{2y}{\varphi(k)\log(y/k)}.
\]
This bounds the number of primes lost to excluded residue classes during
selection \cite{MV73}.
\end{enumerate}

These fifteen named analytic statements, together with the finite certificate
assertion, are the external inputs to the all-odd result presented here.
The range constructions supply the internal prime-pair selections, placements,
completions and numerical inequalities. The full source statements, including
their domains, specify the hypotheses; an asymptotic assertion that primes
are plentiful would not provide the finite constants needed in the walkthrough.

\section*{Acknowledgements}
OpenAI AI systems assisted with drafting and revising the text, developing and
checking the expository structural proofs, reviewing the accompanying
formalization and computational certificates, and preparing the LaTeX source.
The cited analytic theorems are due to their named authors. These assistance
roles do not replace the mathematical arguments, cited construction results,
or exact computations on which the conclusions depend.

\begin{thebibliography}{9}
\bibitem{RS62}
J. Barkley Rosser and Lowell Schoenfeld,
\emph{Approximate formulas for some functions of prime numbers},
Illinois Journal of Mathematics \textbf{6} (1962), 64--94.
Theorems 2--4 and their cited corollaries, pp.~69--70.
\href{https://denisevellachemla.eu/Rosser-Schoenfeld-1962.pdf}{Paper}.
\bibitem{BMOR18}
Michael A. Bennett, Greg Martin, Kevin O'Bryant and Andrew Rechnitzer,
\emph{Explicit bounds for primes in arithmetic progressions},
Illinois Journal of Mathematics \textbf{62} (2018), 427--532.
Theorem 1.9, Propositions 2.1, 2.5, 2.7 and 4.34, and Definitions 2.6 and 3.1.
\href{https://personal.math.ubc.ca/~bennett/BeMaObRe-IJM-2018.pdf}{Paper}.
\bibitem{MV73}
H. L. Montgomery and R. C. Vaughan,
\emph{The large sieve}, Mathematika \textbf{20} (1973), 119--134.
Theorem 2, equation (1.10), p.~121.
\href{https://personal.science.psu.edu/rcv4/personal/Publications/large_sieve.pdf}{Paper}.
\bibitem{Companion}
Mira, \emph{Sude's prime-sum squares: research writeup}, 11 November 2026.
Sections 3.1--3.3 give the mathematical mechanism and source locations;
Section 4 identifies the lower, middle and unbounded construction results;
Sections 6--7 describe the finite certificates and verification.
\href{https://github.com/Mira-Cult/math/blob/d16fd1ee2dd6f8440c6c47b1b7575d0fca8b4241/notes/sudes-first-conjecture/research-writeup.md}{Research writeup and construction source}.
\end{thebibliography}
\end{document}
